\chapter*{INTRODUCCIÓN}
\addcontentsline{toc}{chapter}{\MakeUppercase{INTRODUCCIÓN}}

El término “deuda social” se acuñó por primera vez hace más de 50 años \cite{muir1962socialdebt}, con un enfoque económico para explicar las obligaciones sociales derivadas del intercambio de favores entre personas en una transacción, es decir, entre acreedores y deudores. En otras palabras, cuanto mayor sea el número de relaciones tensas, mayor será la deuda social \cite{caballero2022communitysmells}. Sin embargo, este concepto no se limita a aspectos económicos, ya que su estudio se ha extendido a otros ámbitos, como el desarrollo de software \cite{caballero2022communitysmells}. En este contexto, la deuda social se refiere a las obligaciones pendientes que afectan la calidad, la colaboración y el bienestar en los proyectos de desarrollo \cite{caballero2022communitysmells} \cite{tamburri2013socialdebt}. En esencia, la deuda social encapsula las repercusiones no técnicas derivadas de deci-siones apresuradas o inadecuadas en la gestión de equipos y proyectos, lo que impacta negativamente en la cultura organizacional, la comunicación interna, el bienestar del equipo y, en última instancia, en el éxito del proyecto.
La deuda social puede compararse con un “Leviatán”, un término que evoca la imagen de un monstruo poderoso y temible que crece de manera progresiva y silenciosa, volviéndose cada vez más difícil de ignorar. En el contexto del desarrollo de software, este “Leviatán” representa la acumulación de obligaciones pendientes y tensiones interpersonales que surgen de decisiones apresuradas o inadecuadas en la gestión de equipos y proyectos. Este fenómeno genera un daño colateral que afecta a toda la estructura jerárquica de desarrollo, causando costos y deudas en múltiples aspectos \cite{caballero2022communitysmells}. Uno de los pilares fundamentales de una organización de desarrollo de software, el talento humano, se ve directamente impactado por la deuda social. Las afectaciones se manifiestan sin importar el cargo o función de una persona, evidenciándose principalmente en su rendimiento y calidad de trabajo, los cuales dependen de un estado de bienestar mental y físico, así como de un ambiente laboral justo y motivador \cite{caballero2022communitysmells}. Una mala gestión de estos factores es la causa principal de la alta rotación de personal, la toma de decisiones erróneas y la entrega de productos de baja calidad.
La deuda social surge principalmente cuando se priorizan objetivos a corto plazo, como cumplir con fechas de entrega ajustadas, reducir costos operativos de manera indiscriminada o incorporar requisitos imprevistos en la planificación estratégica inicial, sacrificando las buenas prácticas de gestión y deteriorando las relaciones interpersonales \cite{caballero2021thesis}. Esto conduce a un ambiente de trabajo tóxico, estrés crónico y problemas de salud mental y física, lo que a su vez disminuye la motivación, la productividad y la calidad del trabajo, limitando la capacidad de innovación y creatividad \cite{caballero2022communitysmells}, \cite{tamburri2013socialdebt}, \cite{caballero2021thesis}, \cite{dreesen2021socialdebt}. Para mitigar o gestionar esta deuda, es crucial adoptar un enfoque holístico que promueva la comunicación efectiva, el equilibrio entre trabajo y vida personal, el reconocimiento profesional y un liderazgo que priorice el bienestar del equipo, fomentando un ambiente de trabajo positivo y motivador.
Identificar y comprender la naturaleza y el impacto de la deuda social permite a las organizaciones de software desarrollar estrategias más efectivas para promover un ambiente de trabajo saludable y productivo \cite{dreesen2021socialdebt}. Esto incluye implementar prácticas de gestión que equilibren las demandas operativas con el bienestar de los empleados, fomentando un clima de trabajo sostenible y propicio para la innovación y la creatividad \cite{caballero2022communitysmells}. Las estrategias para gestionar la deuda social deben ser claras y objetivas, permitiendo identificar su impacto en los proyectos. Estas pueden incluir evaluaciones específicas, como pruebas de calidad en procesos de trabajo, encuestas de satisfacción laboral, análisis de rotación de personal y estudios de clima organizacional que reflejen la realidad del equipo de trabajo. Con la información recopilada, es posible gestionar proactivamente la deuda, ajustando políticas de gestión de recursos humanos y estrategias de desarrollo de software ágil. Además, la transparencia en la comunicación y la participación del equipo en la toma de decisiones son vitales para construir un entorno de confianza y colaboración, lo que mejora la moral del equipo y promueve un mayor sentido de pertenencia y responsabilidad colectiva en la calidad del producto final. En este sentido, invertir en la salud organizacional a través de la gestión y mitigación de la deuda social no solo mejora el bienestar de los empleados, sino que también se traduce en una mayor eficiencia operativa, innovación sostenida, participación creativa constante y una ventaja competitiva en el mercado \cite{caballero2022communitysmells}. Por lo tanto, es esencial que las organizaciones reconozcan la deuda social como un aspecto crítico que requiere atención continua y estrategias bien definidas para su gestión efectiva.
Para tener un panorama más amplio del contexto de investigación, se realizó una Revisión Sistemática de la Literatura (RSL) con el fin de recopilar, identificar y priorizar estudios relevantes sobre la gestión de la deuda social en equipos de desarrollo de software ágil. El análisis resultante permitió identificar metodologías y herramientas que intentan mitigar esta deuda, como prácticas de colaboración y comunicación dentro del equipo. Sin embargo, se observó que no existe un gran número de estudios relevantes que aporten significativamente a esta problemática, lo que indica una falta de atención académica y práctica en este ámbito.
Aunque el término “deuda social” se acuñó en el desarrollo de software hace más de una década, no existe un planteamiento claro para gestionarla y mitigarla en equipos de desarrollo de software ágil. Esto resalta la necesidad de desarrollar estrategias efectivas que aborden no solo los aspectos técnicos, sino también las dinámicas interpersonales y organizacionales que influyen en el bienestar del equipo y en la calidad del producto final. La investigación futura debería enfocarse en crear marcos de trabajo que integren la gestión de la deuda social como parte esencial del proceso ágil, promoviendo un entorno más saludable y productivo para los equipos de desarrollo.
Teniendo en cuenta lo anterior y tras analizar la literatura relacionada, surge la siguiente pregunta de investigación: ¿Cómo gestionar la deuda social que surge en equipos de desarrollo de software ágil? Para dar respuesta a este interrogante, se propone adaptar el enfoque ágil para la gestión de proyectos conocido como Scrum, documentando y detallando los elementos esenciales para apoyar la gestión de la deuda social en los proyectos de software. El marco de trabajo ágil propuesto incluirá sesiones de retrospectiva enfocadas en la deuda social, donde el equipo discutirá los desafíos y propondrá soluciones colaborativas. Además, se implementarán sprints dedicados exclusivamente a resolver estos problemas.